% GENERATED FILE. Edit canonical lessons, then run npm run generate:books.
% canonical-source-hash: fnv1a64:f623b9c032b8e7e4
% canonical-lessons: PA-R24-know-think-r4, PA-R24-understand-read-r4, PA-R24-write-take-ask-r4, PA-R24-help-like-r4

\chapter{Old Action Families at R4}
\label{ch:pa-action-r4-bridge}
\hlchaptermodality{24}

\begin{chapteropening}
\textbf{By the end of this chapter:} \emph{I can retrieve the already-taught Punjabi action families at long distance without confusing speech or reading with independent writing.}

The learner completes the action-family R4 bridge through four bounded no-new-language retrieval steps, ending with the distinct help and liking frames.
\end{chapteropening}

\section[Practice]{Know and think return at R4}

\label{lesson:PA-R24-know-think-r4}



\begin{quote}
No new language. Hear \textbf{jāṇnā}, “to know,” and \textbf{sochṇā}, “to think.” Say each meaning before reading on.
\end{quote}

\subsection*{Your turn}
Read \textbf{\pa{ਜਾਣਨਾ} · \pa{ਸੋਚਣਾ}} once, then give the meanings aloud. This is recognition and speech retrieval. It awards no independent Gurmukhi writing evidence.

\begin{tcolorbox}[breakable,colback=teal!4,colframe=teal!35!black,title={Before you move on}]
Without looking, say “to know; to think” in Punjabi.
\end{tcolorbox}
\section[Practice]{Understand and read return at R4}

\label{lesson:PA-R24-understand-read-r4}



\begin{quote}
Say \textbf{samajhṇā} and \textbf{paṛhnā} with the already-taught falling pitch. Give their meanings aloud.
\end{quote}

\subsection*{Your turn}
\textbf{Script recognition only.}

Recognizing \textbf{\pa{ੜ੍ਹ}} here is reading evidence only. Do not copy it, and do not count this lesson as independent writing.

\begin{tcolorbox}[breakable,colback=teal!4,colframe=teal!35!black,title={Before you move on}]
Say “understand; read” once more from meaning.
\end{tcolorbox}
\section[Practice]{Write, take, and ask return at R4}

\label{lesson:PA-R24-write-take-ask-r4}



\begin{quote}
Hear \textbf{likhṇā · laiṇā · puchhṇā}. Give the three meanings aloud before looking at the answer.
\end{quote}

\subsection*{Your turn}
\textbf{Three remembered histories.}

Read \textbf{\pa{ਲਿਖਣਾ} · \pa{ਲੈਣਾ} · \pa{ਪੁੱਛਣਾ}} once. The visible words support recognition; this lesson does not ask for or award independent Gurmukhi writing.

\begin{tcolorbox}[breakable,colback=teal!4,colframe=teal!35!black,title={Before you move on}]
From meaning alone, say “write; take; ask.”
\end{tcolorbox}
\section[Practice]{Help and like return at R4}

\label{lesson:PA-R24-help-like-r4}



\begin{quote}
Say \textbf{madad karnā}, “to help.” Then say the already-known line \textbf{mainū̃ panjābī pasand hai}, “I like Punjabi.”
\end{quote}

\subsection*{Your turn}
\textbf{Liking keeps the liker in the dative.}

Read \textbf{\pa{ਮਦਦ} \pa{ਕਰਨਾ} · \pa{ਮੈਨੂੰ} \pa{ਪੰਜਾਬੀ} \pa{ਪਸੰਦ} \pa{ਹੈ}} once. Romanization records speech retrieval only and earns no Gurmukhi writing credit.

\begin{tcolorbox}[breakable,colback=teal!4,colframe=teal!35!black,title={Before you move on}]
Say “to help” and “I like Punjabi” once more without looking.
\end{tcolorbox}
